\usepackage{fontspec}
\usepackage{polyglossia}
\usepackage{tikz}
\setmainfont{TH Sarabun New}
\newfontfamily\thaifont{TH Sarabun New}
\newfontfamily\thaifonttt{Courier New}
\usepackage{eso-pic}
\usepackage{graphicx}
\setdefaultlanguage{thai}
\setotherlanguages{english}
\usepackage{setspace} 
\usepackage{amsmath}
\usepackage{graphicx}
\usepackage{ragged2e}
\usepackage{svg}
\usepackage{hyperref}
\usepackage{amsthm}
\usepackage{eso-pic}
\usepackage{xcolor}
\usepackage{titlesec}
\usepackage{sectsty}
\usepackage{fancyhdr}
\usepackage{titling}
\usepackage{geometry}
\usepackage{parskip}
\usepackage{booktabs}
\usepackage{threeparttable}
\usepackage{framed}

\geometry{
    top=1in,
    bottom=1in,
    left=1.5in,
    right=1in
}

% การปรับแต่งสารบัญ
\renewcommand{\contentsname}{\centering สารบัญ}
\renewcommand{\listfigurename}{\centering สารบัญภาพ}
\renewcommand{\listtablename}{\centering สารบัญตาราง}
\renewcommand{\chaptername}{\hfill บทที่}
\renewcommand{\figurename}{\textcolor{red}{\textbf{ภาพที่}}}
\renewcommand{\tablename}{\textcolor{red}{\textbf{ตารางที่}}}

% การปรับแต่งชื่อบท Chapter ด้วยกรอบ TikZ
\titleformat{\chapter}[display]
   {\fontsize{28}{28}\color{blue}\bfseries\centering}
   {\begin{tikzpicture}
    \node[rectangle, draw=blue, fill=orange!10, rounded corners, minimum width=15cm, minimum height=1.5cm]
    {\color{blue}\chaptertitlename  \hspace{0.3em}  \thechapter};
    \end{tikzpicture}}{24pt}{\Huge}
\titlespacing*{\chapter}{0pt}{16pt}{0pt}

% หัวข้อย่อย Section ด้วยกรอบ TikZ
\titleformat{\section}
   {\fontsize{16}{18}\color{teal}\bfseries}
   {\begin{tikzpicture}
    \node[rectangle, draw=teal, fill=teal!10, rounded corners, minimum width=8cm, minimum height=1cm]
    {\thesection};
    \end{tikzpicture}}{1em}{}
\titlespacing*{\section}{0pt}{12pt}{8pt}

% การตั้งค่ารูปแบบเนื้อหา
\renewcommand{\normalsize}{\fontsize{16}{16}\selectfont}
\setlength{\parindent}{1.25cm}
\setlength{\parskip}{0.5\baselineskip}

% กรอบสำหรับเนื้อหาด้วย framed
\newenvironment{contentbox}{
    \begin{framed}
    \setlength{\parskip}{0.5\baselineskip}
}{
    \end{framed}
}

% ฟังก์ชันวันที่ภาษาไทย
\newcommand{\thaimonthname}[1]{%
  \ifcase#1
    \or มกราคม\or กุมภาพันธ์\or มีนาคม\or เมษายน\or พฤษภาคม\or มิถุนายน%
    \or กรกฎาคม\or สิงหาคม\or กันยายน\or ตุลาคม\or พฤศจิกายน\or ธันวาคม%
  \fi
}
\newcommand{\thaidate}{%
 \the\day\ \thaimonthname{\month}\ \number\numexpr\year+543\relax
}

% ปรับรูปแบบหน้ากระดาษ
\pagestyle{fancy}
\fancyhf{}
\renewcommand{\headrulewidth}{0pt}
\fancyfoot[C]{\thepage}

% การตั้งค่า Listings
\usepackage{listings}
\lstset{
    basicstyle=\fontsize{12}{14}\ttfamily,
    columns=fullflexible,
    frame=single,
    breaklines=true,
    postbreak=\mbox{\textcolor{red}{$\hookrightarrow$}\space},
    language=Python,
}

% การตั้งค่า Hyperlinks
\hypersetup{
    colorlinks=true,
    linkcolor=blue,
    filecolor=magenta,
    urlcolor=cyan,
    pdftitle={จุลชีววิทยาสำหรับการเกษตร},
    pdfpagemode=FullScreen,
}

% คำสั่งกำหนดเอง
\newcommand{\degree}{\ensuremath{^\circ}}
